\section{Thực nghiệm}
\label{sec:thuc-nghiem}

% Phan thiet ke thi nghiem da co dinh; cac cho [PLACEHOLDER] cho so lieu/cau
% hinh cuoi cung tu terminal thi nghiem (evals/results/).

\subsection{Câu hỏi nghiên cứu}

\begin{itemize}
  \item \textbf{RQ1}: Chiến lược điều phối (ReAct vs.\ Plan-and-Execute) ảnh
        hưởng thế nào đến tỷ lệ thành công, chi phí và độ trễ?
  \item \textbf{RQ2}: Chất lượng mô hình ảnh hưởng thế nào đến các metric trên,
        và có tương tác với chiến lược không (mô hình yếu hơn có hưởng lợi từ
        kế hoạch tường minh)?
  \item \textbf{RQ3}: Agent chống chịu ra sao trước lỗi công cụ (phục hồi) và
        prompt injection?
\end{itemize}

\subsection{Thiết kế}

Ma trận $2 \times 2$: \{ReAct, Plan-and-Execute\} $\times$ \{%
% PLACEHOLDER: ten 2 model cuoi cung, vd gemini-3.1-flash-lite va gpt-4o-mini
\emph{model A}, \emph{model B}\}, mỗi cấu hình chạy toàn bộ 37 case với
$N \geq 3$ lượt/case (LLM không tất định). Bộ case và prompt được đóng băng
trước khi chạy; mỗi lượt chạy trên một DB sạch, cô lập hoàn toàn.

Metric thu thập tự động từ trace (xem Mục~\ref{sec:kien-truc}): tỷ lệ thành
công, độ chính xác chọn tuyến, độ chính xác chọn công cụ, tỷ lệ phục hồi,
tỷ lệ đạt LLM-as-judge, số bước trung bình, số lệnh gọi LLM trung bình,
token + chi phí trung bình mỗi tác vụ, thời gian chạy trung bình.

\subsection{Môi trường}

% PLACEHOLDER: dien phien ban/cau hinh cu the khi chot so lieu.
Máy chạy: [PLACEHOLDER --- CPU/RAM/OS]. Python [3.x], mã nguồn tại commit
[\texttt{abcdef0}]. Model chấm judge: [PLACEHOLDER --- phải khác model agent].
Giới hạn an toàn: tối đa 8 bước/tác vụ, timeout công cụ 15\,s, timeout tác vụ
180\,s, retry công cụ 2 lần.
